% D8 — maquette du livret de l'élève, en A5 : HUIT pages exactement, pour
% tenir sur deux feuilles A4 pliées (voir D8-livret-a-imprimer.tex).
%   1 couverture · 2-3 parcours · 4-5 savoir-faire · 6-7 carnet d'écoute
%   8 attestation
% Savoir-faire repris de bir/referentiel.md.
\documentclass[a5]{BIRdoc}
\BIRdocument{Livret de l'élève}{}
\newcommand{\savoir}[2]{#1 & #2 & \BIRcase \\}
\newenvironment{savoirfaire}{%
  \small\tabularx{\linewidth}{P{10mm}LP{11mm}}
  \BIRentete \BIRth{Séance} & \BIRth{Je sais…} & \BIRth{Validé} \\}
  {\bottomrule\endtabularx\par}
\newenvironment{parcours}{%
  \small\tabularx{\linewidth}{|P{6mm}|L|P{24mm}|P{22mm}|}
  \hline\BIRentete \BIRth{Nº} & \BIRth{Séance} & \BIRth{Date} & \BIRth{Visa} \\ \hline}
  {\endtabularx\par}
\newcommand{\seance}[2]{\rule{0pt}{9.2mm}#1 & #2 & & \\ \hline}
\newenvironment{carnet}{%
  \footnotesize\tabularx{\linewidth}{|P{13mm}|P{10mm}|P{17mm}|P{22mm}|L|}
  \hline\BIRentete \BIRth{Date} & \BIRth{Heure} & \BIRth{Fréquence} & \BIRth{Indicatif, station} & \BIRth{Ce que j'ai entendu} \\ \hline}
  {\endtabularx\par}
\newcommand{\ecoute}{\rule{0pt}{7mm} & & & & \\ \hline}
\begin{document}

% ---------------------------------------------------------------- 1 couverture
\thispagestyle{empty}
\vspace*{8mm}
{\sffamily\bfseries\fontsize{27}{32}\selectfont\color{BIRbleu}Brevet d'Initiation\\ à la Radio\par}
\vspace{2mm}
{\color{BIRrouge}\rule{40mm}{1.2pt}\par}
\vspace{3mm}
{\sffamily\LARGE Mon livret\par}
\vspace{18mm}
Nom : \BIRligne[8.5cm]\par\vspace{7mm}
Prénom : \BIRligne[7.9cm]\par\vspace{7mm}
Établissement : \BIRligne[6.8cm]\par\vspace{7mm}
Année scolaire : \BIRligne[3.5cm]\par\vspace{7mm}
Mon encadrant : \BIRligne[6.7cm]\par\vspace{7mm}
Son indicatif : \BIRligne[3.5cm]\par
\vfill
\begin{BIRencadre}{À quoi sert ce livret ?}
Il me suit pendant les 24 séances. J'y note ma présence, ce que je sais
faire, et ce que j'entends sur les ondes. À la fin du parcours, il reçoit
mon attestation.
\end{BIRencadre}

% ---------------------------------------------------------------- 2-3 parcours
\newpage
\BIRtitre{Mon parcours}{Je note la date, mon encadrant signe}
\begin{parcours}
\seance{1}{La radio autour de nous}
\seance{2}{Se présenter à l'antenne}
\seance{3}{Un contact radio}
\seance{4}{Tension et courant}
\seance{5}{La résistance et la loi d'Ohm}
\seance{6}{Puissance et énergie}
\seance{7}{Du courant alternatif à l'onde}
\seance{8}{Le spectre et les bandes}
\seance{9}{La propagation à vue}
\seance{10}{Les ondes courtes et l'ionosphère}
\seance{11}{Moduler : transporter la voix}
\seance{12}{Le Morse}
\end{parcours}

\newpage
\begin{parcours}
\seance{13}{Les modes numériques et l'image}
\seance{14}{Les antennes : principes}
\seance{15}{Construire son antenne}
\seance{16}{Câbles, connecteurs et réglage}
\seance{17}{Le récepteur}
\seance{18}{L'émetteur}
\seance{19}{L'atelier de soudure}
\seance{20}{La radio en action}
\seance{21}{Préparer son contact}
\seance{22}{Au micro}
\seance{23}{Les règles, et après ?}
\seance{24}{L'examen du BIR}
\end{parcours}

% ---------------------------------------------------------------- 4-5 savoir-faire
\newpage
\BIRtitre{Ce que je sais faire}{Mon encadrant coche quand il m'a vu réussir}

\section*{A --- Découvrir la radio}
\begin{savoirfaire}
\savoir{1}{régler un récepteur sur une station et la décrire}
\savoir{20}{trouver mon locator ; localiser un émetteur caché}
\end{savoirfaire}

\section*{B --- Électricité et électronique}
\begin{savoirfaire}
\savoir{4}{câbler un circuit simple ; mesurer une tension et un courant}
\savoir{5}{calculer avec la loi d'Ohm et vérifier par la mesure}
\savoir{6}{calculer une puissance et une autonomie}
\savoir{17}{régler fréquence, volume, mode et squelch d'un poste}
\savoir{18}{lire une puissance sur un wattmètre}
\savoir{19}{réaliser une soudure propre ; suivre un schéma d'implantation}
\end{savoirfaire}

\section*{C --- Ondes, spectre et propagation}
\begin{savoirfaire}
\savoir{7}{calculer une longueur d'onde ; lire une période sur un oscillogramme}
\savoir{8}{lire un spectre et une chute d'eau ; dresser une carte du spectre}
\savoir{9}{mesurer une portée et la reporter sur un plan}
\savoir{10}{identifier une station lointaine et la placer sur une carte}
\end{savoirfaire}

\newpage
\section*{D --- Modulations, antennes et station}
\begin{savoirfaire}
\savoir{11}{reconnaître un mode à l'oreille et à l'écran}
\savoir{12}{émettre et recevoir un mot court en Morse}
\savoir{13}{décoder une transmission avec un logiciel}
\savoir{14}{calculer la longueur d'un dipôle}
\savoir{15}{assembler une antenne ; constater sa directivité}
\savoir{16}{mesurer et améliorer le réglage d'une antenne}
\end{savoirfaire}

\section*{E --- Trafic, règles et sécurité}
\begin{savoirfaire}
\savoir{2}{épeler et noter un indicatif sans erreur}
\savoir{3}{tenir un carnet de trafic}
\savoir{21}{dérouler un contact complet, script en main}
\savoir{22}{tenir ma place dans un vrai contact, aux côtés du radioamateur}
\end{savoirfaire}

\vfill
\begin{BIRencadre}{Mes remarques}
\vspace{30mm}
\end{BIRencadre}

% ---------------------------------------------------------------- 6-7 carnet
\newpage
\BIRtitre{Mon carnet d'écoute}{Une ligne par station entendue}
\begin{carnet}
\ecoute\ecoute\ecoute\ecoute\ecoute\ecoute\ecoute\ecoute\ecoute\ecoute
\ecoute\ecoute\ecoute\ecoute\ecoute\ecoute
\end{carnet}

\newpage
\begin{carnet}
\ecoute\ecoute\ecoute\ecoute\ecoute\ecoute\ecoute\ecoute\ecoute\ecoute
\ecoute\ecoute\ecoute\ecoute\ecoute\ecoute\ecoute\ecoute\ecoute
\end{carnet}

% ---------------------------------------------------------------- 8 attestation
\newpage
\thispagestyle{empty}
\vspace*{6mm}
\begin{center}
{\sffamily\bfseries\fontsize{22}{27}\selectfont\color{BIRbleu}Brevet d'Initiation\\ à la Radio\par}
\vspace{3mm}
{\color{BIRrouge}\rule{40mm}{1.2pt}\par}
\vspace{5mm}
{\sffamily\LARGE Attestation\par}
\vspace{14mm}
Il est attesté que\par\vspace{8mm}
\BIRligne[9cm]\par\vspace{8mm}
a suivi les 24 séances de préparation\\ au Brevet d'Initiation à la Radio\par
\vspace{5mm}
et a réussi l'examen du \BIRligne[3.5cm]\par
\vspace{3mm}
avec \BIRligne[1.5cm] bonnes réponses sur 40.\par
\vspace{18mm}
À \BIRligne[4cm], le \BIRligne[3cm]\par
\vspace{12mm}
\begin{tabularx}{\linewidth}{LL}
L'encadrant & Le responsable de l'établissement \\[1mm]
{\small nom, indicatif, signature} & {\small nom, signature} \\
\end{tabularx}
\end{center}

\end{document}
